\section{Toán học, công cụ để làm khoa học}

\ 

Chắc các bạn đọc từng nghe ngạn ngữ Anh: ``Một bức hình thì giá trị bằng cả ngàn câu từ'', và thực tế, chúng ta không thể phủ nhận tính súc tích và dễ hiểu của một tranh minh họa được đặt đúng vị trí. Các phương trình toán học cũng như vậy. Hãy so sánh hai cách biểu diễn, cho cùng một định luật quen thuộc:

\begin{table}[H]
    \centering
    \caption{So sánh định luật biểu diễn bằng lời và bằng phương trình toán}
    \begin{tblr}{
        colspec = {Q[m,l,8cm] Q[m,c,5cm]}, % Q là cột thông minh
        % m: căn giữa dọc, l: căn trái ngang, c: căn giữa ngang
        hlines, vlines, % Kẻ bảng cực nhanh
        rows = {1.5ex}, % Độ cao dòng
        hline{2} = {1.5pt}, 
        vline{2} = {1.5pt}
    }
        \textbf{Biểu diễn bằng lời} & \textbf{Biểu diễn bằng phương trình} \\
        Gia tốc của một chất điểm thì tỉ lệ thuận với lực tác động lên vật và tỉ lệ nghịch với khối lượng của vật. & 
        $\vec{a} = \frac{\vec{F}}{m}$. \\
    \end{tblr}
\end{table}